% ============================================================
% macros_annotations.tex — \ptFleche et style de flèches annotées
%
% \ptFleche{nom}{contenu} : place CONTENU normalement (dans une équation,
% typiquement) tout en le nommant "nom" pour pouvoir y tracer une flèche
% ensuite depuis un tikzpicture[overlay] séparé. Contrairement à
% \tikzmark{nom} (repère invisible AUTOUR d'un texte existant),
% \ptFleche ENVELOPPE le contenu -- pratique pour annoter un terme d'une
% équation (\begin{eqnarray}...\ptFleche{LL}{$\operator{H}$}...\end{eqnarray})
% sans avoir à répéter le terme deux fois.
%
% ⚠️ Comme pour overlay_tikz : place le tikzpicture qui DESSINE les
% flèches juste après le bloc qui contient les \ptFleche référencés
% (même question/section), jamais en fin d'exercice -- sinon risque de
% décalage de page (bug réel rencontré, voir data/SCHEMA.md).
% ============================================================
\usetikzlibrary{arrows}   % pour le style de pointe "Computer Modern Rightarrow"

% Notation opérateur (physique quantique) : \operator{H} -> Ĥ
\newcommand{\operator}[1]{\hat{#1}}

\newcommand{\ptFleche}[2]{\tikz[remember picture, baseline=(#1.base)]{\node[inner sep=0pt, outer sep=0pt] (#1) {#2};}}


% ------------------------------------------------------------------
% 1. CONTRÔLE PRÉCIS DES MARGES (geometry)
% ------------------------------------------------------------------

\geometry{
  a4paper,
  top=1.8cm,         % Marge supérieure
  bottom=1.8cm,      % Marge inférieure
  left=1.5cm,        % Marge gauche
  right=1.5cm,       % Marge droite
  headheight=14pt,   % Hauteur réservée à l'en-tête (évite les warnings fancyhdr)
  headsep=0.6cm,     % Distance entre l'en-tête et le corps du texte
  footskip=0.8cm     % Distance entre le corps du texte et le pied de page
}

% ------------------------------------------------------------------
% 2. COULEURS DES LIENS (\ref, citations, URL)
% ------------------------------------------------------------------
\definecolor{MyLinkColor}{HTML}{1565C0}  % Couleur pour les \ref, TOC, etc.
\definecolor{MyCiteColor}{HTML}{2E7D32}  % Couleur pour les citations biblio
\definecolor{MyUrlColor}{HTML}{B71C1C}   % Couleur pour les hyperliens http/https

\hypersetup{
  colorlinks=true,         % Active la coloration des liens (au lieu des cadres)
  linkcolor=MyLinkColor,   % Appliqué aux \ref{}, \pageref{}, sommaire, etc.
  citecolor=MyCiteColor,   % Appliqué aux \cite{}
  urlcolor=MyUrlColor,     % Appliqué aux \url{} et \href{}
  pdfborder={0 0 0}        % Supprime tout encadré résiduel
}

% ------------------------------------------------------------------
% 3. EN-TÊTE ET PIED DE PAGE AVEC LIGNES PERSONNALISÉES (fancyhdr)
% ------------------------------------------------------------------

% 2. CONFIGURATION CRUCIALES : Recalculer la largeur pour ajuster les traits
\fancyhfoffset[L,R]{0pt} % Aligne exactement les traits sur le texte
\setlength{\headwidth}{\textwidth} % Force la largeur du trait à suivre les marges

\usepackage{fancyhdr}
\pagestyle{fancy}
\fancyhf{} % Réinitialise les en-têtes et pieds de page

% --- Couleurs et épaisseurs des lignes
\definecolor{HeaderRuleColor}{HTML}{4A7AA4} % Couleur du trait du haut
\definecolor{FooterRuleColor}{HTML}{7A7A7A} % Couleur du trait du bas

\renewcommand{\headrulewidth}{1.5pt} % Épaisseur du trait en haut
\renewcommand{\footrulewidth}{0.8pt} % Épaisseur du trait en bas

% Coloration dynamique des filet/lignes supérieur et inférieur
\renewcommand{\headrule}{{\color{HeaderRuleColor}\hrule width\headwidth height\headrulewidth \vspace{-\headrulewidth}}}
\renewcommand{\footrule}{{\color{FooterRuleColor}\hrule width\headwidth height\footrulewidth}}

% --- Contenu de l'en-tête (Header)
\fancyhead[L]{\small\bfseries\color{woodNoyer} Mon Cours / Mon TD}
\fancyhead[R]{\small\itshape\color{grayPierre} Année 2025–2026}

% --- Contenu du pied de page (Footer)
\fancyfoot[L]{\small\color{grayBrume} Université / Établissement}
\fancyfoot[C]{\small\color{grayPierre} \thepage} % Numéro de page
\fancyfoot[R]{\small\color{grayBrume} Module : Physique des Ondes}

% --- Appliquer également le style aux pages de garde / début de chapitre
\fancypagestyle{plain}{
  \fancyhf{}
  \renewcommand{\headrulewidth}{1.5pt} % 0pt Pas de ligne en haut sur la page de titre
  \renewcommand{\footrulewidth}{0.8pt}
  \renewcommand{\footrule}{{\color{FooterRuleColor}\hrule width\headwidth height\footrulewidth}}
  \fancyfoot[L]{\small\color{grayBrume} Université / Établissement}
  \fancyfoot[C]{\small\color{grayPierre} \thepage} % Numéro de page
  \fancyfoot[R]{\small\color{grayBrume} Module : Physique des Ondes}
}
